\documentclass[tikz,border=10pt]{standalone}
\usepackage{ctex}
\usepackage{amsmath}
\usetikzlibrary{calc, arrows.meta, angles, quotes}

\begin{document}
% geo-p9-05-1: 3 类空间角的几何意义对比（线线 / 线面 / 二面角）
\begin{tikzpicture}[
  x={(1cm,0cm)},
  y={(0.40cm,0.32cm)},
  z={(0cm,1cm)},
  thick,
  >=Stealth,
  line cap=round, line join=round,
  label style/.style={font=\small},
  angle color/.style={draw=red!70!black, fill=red!15},
]

%% ============================================================
%% 栏 1：线线角（两异面直线夹角，取锐角）
%% ============================================================
\begin{scope}[xshift=0cm]
  % 浅色参考平面
  \fill[blue!8, draw=blue!30, thin]
    (-1.0, -1.0, 0) -- (3.5, -1.0, 0) -- (3.5, 2.5, 0) -- (-1.0, 2.5, 0) -- cycle;

  % 直线 l1（蓝，在平面内方向）
  \coordinate (L1A) at (-0.5, 0.0, 0);
  \coordinate (L1B) at (3.0, 0.0, 0);
  \draw[blue!70!black, thick, ->] (L1A) -- (L1B) node[right, font=\small] {$l_1$};

  % 直线 l2（绿，空间中另一方向，起点偏上）
  \coordinate (L2A) at (0.5, 2.0, 1.2);
  \coordinate (L2B) at (3.0, 0.8, 0.3);
  \draw[green!60!black, thick, ->] (L2A) -- (L2B) node[right, font=\small] {$l_2$};

  % 平移 l2 到与 l1 相交（O 点），画夹角
  \coordinate (O1) at (1.5, 0.0, 0);
  % l2 方向向量 = L2B - L2A，归一化后展示
  % 简化：直接画两条共点的方向箭头表示夹角
  \coordinate (V2) at (2.5, 0.5, -0.5);  % l2 的平移副本方向端点
  \draw[green!60!black, dashed, ->] (O1) -- (V2) node[right, font=\small] {$l_2'$（平移）};

  % 标注夹角 θ（取锐角）
  \draw[red!70!black, thick] (O1) -- ++(1.0, 0, 0) coordinate (T1R);
  \draw[red!80, thick, ->] (O1) -- ($(O1)!1.0!(V2)$) coordinate (T1U);
  % 角弧
  \draw[red!60, thick] ($(O1)!0.35!(T1R)$) arc[start angle=0, end angle=23, radius=0.55cm]
    node[midway, right=2pt, font=\small, red!80!black] {$\theta$};

  \node[font=\small\bfseries] at (1.2, -1.3, 0) {(1) 线线角};
  \node[font=\small] at (1.2, -1.8, 0) {$\cos\theta=\left|\dfrac{\vec{a}\cdot\vec{b}}{|\vec{a}||\vec{b}|}\right|$，取绝对值};
  \node[font=\small, align=center] at (1.2, -2.5, 0) {异面线夹角 $\theta\in(0,\pi/2]$};
\end{scope}

%% ============================================================
%% 栏 2：线面角（直线与平面所成角）
%% ============================================================
\begin{scope}[xshift=5.5cm]
  % 平面（灰绿）
  \fill[green!10, draw=green!50!black, thin]
    (-1.0, -1.0, 0) -- (3.5, -1.0, 0) -- (3.5, 2.5, 0) -- (-1.0, 2.5, 0) -- cycle;
  \node[green!50!black, font=\small] at (3.2, 2.2, 0) {$\alpha$};

  % 法向量 n（淡蓝箭头，垂直于平面）
  \coordinate (O2) at (1.5, 1.0, 0);
  \draw[blue!60, ->, dashed] (O2) -- ++(0, 0, 2.0) node[above, font=\small, blue!70] {$\vec{n}$};

  % 直线 l（红，穿过平面）
  \coordinate (lA) at (0.2, -0.5, -0.8);
  \coordinate (lB) at (2.8, 2.5, 1.5);
  \draw[red!70!black, thick, ->] (lA) -- (lB) node[above right, font=\small] {$l$};

  % 射影（l 在平面内的投影，虚线橙色）
  \coordinate (lFoot) at (1.55, 1.1, 0);   % l 与平面交点
  \coordinate (lProj) at (2.8, 2.5, 0);    % 投影方向端点
  \draw[orange!80!black, dashed, thick, ->] (lFoot) -- (lProj) node[right, font=\small] {射影};

  % 线面角 φ
  \draw[red!60, thick] (lFoot) -- ++(0.5, 0.4, 0.0) coordinate (ProjEnd2);
  \draw[red!70!black, thick, ->] (lFoot) -- ++(0.5, 0.4, 0.6) coordinate (LineEnd2);
  \draw[red!50, thick] ($(lFoot)!0.5!(ProjEnd2)$) arc[start angle=38, end angle=55, radius=0.6cm]
    node[midway, right=1pt, font=\small, red!80!black] {$\varphi$};

  \node[font=\small\bfseries] at (1.2, -1.3, 0) {(2) 线面角};
  \node[font=\small] at (1.2, -1.9, 0) {$\sin\varphi=\left|\dfrac{\vec{l}\cdot\vec{n}}{|\vec{l}||\vec{n}|}\right|$};
  \node[font=\small, align=center] at (1.2, -2.5, 0) {线面角 $\varphi\in[0,\pi/2]$，取绝对值};
\end{scope}

%% ============================================================
%% 栏 3：二面角（两半平面夹角）
%% ============================================================
\begin{scope}[xshift=11cm]
  % 棱 l（黑粗线）
  \coordinate (EdgeA) at (1.5, -0.5, 0);
  \coordinate (EdgeB) at (1.5,  2.5, 0);
  \draw[black, very thick] (EdgeA) -- (EdgeB) node[right, font=\small] {棱 $l$};

  % 半平面 α（蓝）
  \fill[blue!15, opacity=0.8]
    (EdgeA) -- (EdgeB) -- (-0.5, 2.5, 0) -- (-0.5, -0.5, 0) -- cycle;
  \draw[blue!60, thin]
    (EdgeA) -- (EdgeB);
  \node[blue!70, font=\small] at (0.2, 1.5, 0.3) {$\alpha$};

  % 半平面 β（绿，向上倾斜）
  \fill[green!15, opacity=0.8]
    (EdgeA) -- (EdgeB) -- (3.5, 2.5, 1.5) -- (3.5, -0.5, 1.5) -- cycle;
  \node[green!60!black, font=\small] at (3.0, 1.5, 1.3) {$\beta$};

  % 棱上一点 O，各作垂直棱的射线
  \coordinate (O3) at (1.5, 1.0, 0);
  \fill (O3) circle [radius=1.5pt];

  % 在 α 中垂直棱的射线 OA（向左）
  \draw[blue!80, thick, ->] (O3) -- ++(-1.2, 0, 0) node[left, font=\small] {$OA$};

  % 在 β 中垂直棱的射线 OB（向右斜上）
  \draw[green!70!black, thick, ->] (O3) -- ++(1.2, 0, 0.9) node[right, font=\small] {$OB$};

  % 法向量 n1（蓝，垂直 α 面，朝上）
  \draw[blue!50, dashed, ->] (O3) -- ++(0, 0, 1.5) node[above, font=\small, blue] {$\vec{n_1}$};

  % 法向量 n2（绿，垂直 β 面，朝上偏右）
  \draw[green!60!black, dashed, ->] (O3) -- ++(-0.7, 0, 1.5) node[above left, font=\small, green!70!black] {$\vec{n_2}$};

  % 二面角弧
  \draw[red!60, thick] (-0.5, 1.0, 0) arc[start angle=180, end angle=143, radius=1.0cm]
    node[midway, below left, font=\small, red!80!black] {$\psi$};

  \node[font=\small\bfseries] at (1.5, -1.3, 0) {(3) 二面角};
  \node[font=\small] at (1.5, -1.9, 0) {$\cos\psi=\pm\dfrac{\vec{n_1}\cdot\vec{n_2}}{|\vec{n_1}||\vec{n_2}|}$};
  \node[font=\small, align=center] at (1.5, -2.5, 0) {不取绝对值，看图判正负};
\end{scope}

\end{tikzpicture}
\end{document}
